\question\TFQuestion{T}{One factor that can cause a population to grow smaller is an increasing death rate.}
\question\TFQuestion{T}{One factor that can cause a population to grow larger is an decreasing death rate.}
\question\TFQuestion{T}{One factor that can cause a population to grow larger is an increasing birth rate.}
\question\TFQuestion{T}{One factor that can cause a population to grow smaller is an decreasing immigration rate.}
\question\TFQuestion{T}{One factor that can cause a population to grow smaller is an increasing emigration rate.}

\question\TFQuestion{F}{One factor that can cause a population to grow larger is an increasing death rate.}
\question\TFQuestion{F}{One factor that can cause a population to grow smaller is an decreasing death rate.}
\question\TFQuestion{F}{One factor that can cause a population to grow smaller is an increasing birth rate.}
\question\TFQuestion{F}{One factor that can cause a population to grow larger is an decreasing immigration rate.}
\question\TFQuestion{F}{One factor that can cause a population to grow larger is an increasing emigration rate.}


\question\TFQuestion{F}{A population will always grow smaller if emigration is greater than birth rate.}
\question\TFQuestion{T}{A population may grow smaller if emigration is greater than birth rate.}

\question\TFQuestion{F}{A population will always grow larger if immigration is greater than death rate.}
\question\TFQuestion{T}{A population may grow larger if immigration is greater than death rate.}

\question\TFQuestion{F}{A population will always grow smaller if immigration is greater than death rate.}
\question\TFQuestion{T}{A population may grow smaller if death rate is greater than immigration rate.}


\question\TFQuestion{F}{Density-independent forces are most likely to affect population size when \textit{N} is less than \textit{K}.}
\question\TFQuestion{F}{Density-independent forces are most likely to affect population size when \textit{N} is greater than \textit{K}.}
\question\TFQuestion{F}{Density-dependent forces are most likely to affect population size when \textit{N} is less than \textit{K}.}
\question\TFQuestion{T}{Density-dependent forces are most likely to affect population size when \textit{N} is greater than \textit{K}.}


\question\TFQuestion{F}{Negative feedback most often occurs in a population when \textit{N} is much less than \textit{K}.}
\question\TFQuestion{T}{Negative feedback most often occurs in a population when \textit{N} is much larger than \textit{K}.}
\question\TFQuestion{T}{The carrying capacity of a habitat may change as environmental conditions change.}
\question\TFQuestion{F}{The carrying capacity of source habitats decreases when the carrying capacity of sink patches increases.}


\question\TFQuestion{F}{Most organisms grow under idealized conditions because they are adapted to their environment.}
\question\TFQuestion{T}{The intrinsic rate of population growth does not include immigration or emigration.}
\question\TFQuestion{F}{The intrinsic rate of population growth includes immigration and emigration.}
\question\TFQuestion{F}{Populations always grow fastest when \emph{N} is less than \emph{K}.}

\question\TFQuestion{T}{Mutualism can be thought of as two species that take advantage of each other ecologically over the long term.}
\question\TFQuestion{T}{Richness measures only the number of species in a community. }
\question\TFQuestion{T}{Evenness measures the distribution of relative abundance across all species in a community.}
\question\TFQuestion{F}{Evenness measures the distribution of relative abundance across all species in a population.}

\question\TFQuestion{T}{The brightly colored pattern of poison dart frogs is an example of aposomatic coloration. }
\question\TFQuestion{T}{The brightly colored pattern of bees and wasps is an example of aposomatic coloration. }
\question\TFQuestion{F}{The brightly colored pattern of harmless coral reef fishes is an example of aposomatic coloration. }
\question\TFQuestion{T}{Competitive exclusion occurs when one species outcompetes another species that uses the realized same niche.}
\question\TFQuestion{T}{Disease killed off most of the sea urchins around Jamaican coral reefs.  The number of species on the coral reef decreased dramatically after the sea urchins died. The sea urchins may have been a keystone species.}

\question\TFQuestion{T}{The carrying capacity of a habitat may change as environmental conditions change.}

\question\TFQuestion{F}{Competitive exclusion will occur when two species use even one identical resource out of many resources.}
\question\TFQuestion{F}{An insect that has evolved to resemble a plant twig will probably be able to avoid parasitism.}
\question\TFQuestion{T}{An insect that has evolved to resemble a plant twig will have a better chance of avoiding predation.}

\question\TFQuestion{F}{Random dispersion occurs when all individuals in a population are evenly spaced.}


\question\TFQuestion{F}{A molecular clock that has been calculated for one gene can be used safely for another gene.}
\question\TFQuestion{T}{A molecular clock that has been calculated for one gene cannot be used safely for another gene.}


\question\TFQuestion{F}{The evolution of large brain size then allowed \textit{Homo} to evolve bipedalism.}

% Meh
\question\TFQuestion{T}{Population is the lowest hierarchical level for the biosphere category.}
\question\TFQuestion{F}{Species is the lowest hierarchical level for the biosphere category.}
\question\TFQuestion{F}{Species is the lowest hierarchical level for the biosphere category.}


\question\TFQuestion{T}{The fossil record can be used to calibrate a molecular clock.}
